\documentclass[11pt,a4paper]{article}
\usepackage[T1]{fontenc}\usepackage[utf8]{inputenc}\usepackage{lmodern}\usepackage{microtype}
\usepackage[a4paper,margin=2.25cm]{geometry}\usepackage{longtable,booktabs,array,calc}\usepackage{fancyvrb}
\usepackage{xcolor}\usepackage{hyperref}\usepackage{xurl}\usepackage{fancyhdr}\usepackage{enumitem}\usepackage{parskip}\usepackage{titlesec}
\definecolor{ddtadark}{RGB}{30,38,48}\definecolor{ddtablue}{RGB}{38,82,126}
\hypersetup{colorlinks=true,linkcolor=ddtablue,urlcolor=ddtablue,citecolor=ddtablue}
\pagestyle{fancy}\fancyhf{}\lhead{Documentation-Driven Threat Analysis}\rhead{BA3-T2 pressure test}\cfoot{\thepage}\setlength{\headheight}{14pt}
\setlist[itemize]{leftmargin=1.5em,itemsep=0.2em,topsep=0.25em}\setlist[enumerate]{leftmargin=1.7em,itemsep=0.2em,topsep=0.25em}
\titleformat{\section}{\Large\bfseries\color{ddtadark}}{}{0em}{}\titleformat{\subsection}{\large\bfseries\color{ddtadark}}{}{0em}{}\titleformat{\subsubsection}{\normalsize\bfseries\color{ddtadark}}{}{0em}{}
\providecommand{\tightlist}{\setlength{\itemsep}{0pt}\setlength{\parskip}{0pt}}\setlength{\emergencystretch}{4em}\hbadness=10000\hfuzz=20pt\sloppy
\begin{document}
\section{DDTA BA3-T2 cross-baseline identity, staleness and lifecycle
pressure test -
R1}\label{ddta-ba3-t2-cross-baseline-identity-staleness-and-lifecycle-pressure-test---r1}

\textbf{Status:} COMPLETED / PROVISIONAL PASS WITH IDENTITY-LIFECYCLE
REFINEMENT

\textbf{Repository baseline reviewed:}
\texttt{52864d2ce177abdd694436306d8152db688effa0}

\textbf{Executed microstep:}
\texttt{BA3-T2\ -\ cross-baseline\ identity,\ staleness\ and\ lifecycle\ pressure\ test}

\textbf{Scope discipline:} this test does not close BA3, does not
redesign BA2, does not define BA4 projections, AnalysisRecord/Common
Finding, STRIDE overlays or ThreatForge implementation state.

\subsection{1. Question under test}\label{question-under-test}

BA3-T1 proved that every BA identity needs baseline-scoped
origin/provenance. T2 asks what happens when the governed baseline
changes:

\begin{quote}
Which BA meanings remain the same identities, which must be reviewed,
which are replaced, and which disappear?
\end{quote}

The test deliberately uses concrete governed mutations instead of
inventing a generic lifecycle state machine.

\subsection{2. Candidate attacked}\label{candidate-attacked}

T1 left these questions open:

\begin{itemize}
\tightlist
\item
  cross-baseline identity/equivalence;
\item
  accepted/rejected/stale/superseded/retired semantics;
\item
  diagnostic resolution after source correction.
\end{itemize}

The pressure hypothesis was intentionally minimal:

\begin{verbatim}
H1 same source ID -> same BA identity
H2 one lifecycle rule can serve BAReferent and BAProposition
H3 source revise can be copied as BA revise
H4 originState can carry acceptance/freshness meaning
\end{verbatim}

All four hypotheses were attacked by the mutation corpus.

\subsection{3. Test A - M1 Ethernet to
Wi-Fi}\label{test-a---m1-ethernet-to-wi-fi}

\subsubsection{3.1 Governed mutation}\label{governed-mutation}

The facial-access corpus changes only \texttt{D-3.5}:

\begin{verbatim}
Ethernet available -> Wi-Fi available
\end{verbatim}

Transport remains externally consumed; camera and recognition processor
remain separated; FR-3.4 and its specialized constraints survive after
review.

\subsubsection{3.2 Identity pressure}\label{identity-pressure}

The abstract \texttt{LocalConnectivity} meaning survives. The concrete
realization does not.

If same document identity implied same BA proposition identity, the
assertion:

\begin{verbatim}
realize(LocalConnectivity, WiredEthernet)
\end{verbatim}

would have to mutate in place into:

\begin{verbatim}
realize(LocalConnectivity, WiFi)
\end{verbatim}

That would make proposition identity denote two different normalized
assertions over history.

\subsubsection{3.3 Result}\label{result}

\begin{verbatim}
LocalConnectivity referent                         RETAIN
WiredEthernet referent                            RETIRE if unused elsewhere
WiFi referent                                     INTRODUCE
old realize proposition                           REPLACE / SUPERSEDED
transfer/correlation/security propositions        RETAIN after review
same D-3.5 source identity -> same BA proposition REJECTED
\end{verbatim}

\textbf{Finding A:} source-document continuity and BA assertion
continuity are independent decisions.

\subsection{4. Test B - M2 external to project-owned
transport}\label{test-b---m2-external-to-project-owned-transport}

\subsubsection{4.1 Governed mutation}\label{governed-mutation-1}

\texttt{D-3.4} changes from consumed external transport to project
ownership/management.

\subsubsection{4.2 Polarity pressure}\label{polarity-pressure}

BA2 represents the V0 boundary as a negative responsibility assertion.
Project ownership changes that truth value.

Keeping one proposition identity while flipping polarity would violate
assertion identity:

\begin{verbatim}
assignResponsibility[negative](project, transport, ownership/management)
\end{verbatim}

becomes:

\begin{verbatim}
assignResponsibility[affirmative](project, transport, ownership/management)
\end{verbatim}

\subsubsection{4.3 Result}\label{result-1}

\begin{verbatim}
abstract transport/connectivity referent          RETAIN if meaning survives
negative responsibility proposition               REPLACE / SUPERSEDED
affirmative responsibility proposition            INTRODUCE as successor
external consumeService proposition               RETIRE if consumption ceases
transfer + specialized constraints                RETAIN after review
new project-owned transport branch                INTRODUCE as needed
\end{verbatim}

\textbf{Finding B:} changing responsibility relations does not
automatically split the referent whose responsibility changed.

\subsection{5. Test C - M3 remote to local
recognition}\label{test-c---m3-remote-to-local-recognition}

\subsubsection{5.1 Governed mutation}\label{governed-mutation-2}

Recognition moves onto the camera. The remote RecognitionCapture
transfer disappears.

\subsubsection{5.2 Retire versus replace
pressure}\label{retire-versus-replace-pressure}

A bad lifecycle rule would transform the old transfer proposition into
any new local processing behavior merely because both concern
recognition.

That would destroy the proposition's identity semantics.

\subsubsection{5.3 Result}\label{result-2}

\begin{verbatim}
CameraSubsystem referent                           RETAIN
RecognitionCapture referent
  RETAIN if acquisition meaning survives
RecognitionProcessor                               REVIEW / possibly RETIRE
remote delivery behavior referent                  RETIRE
remote transfer proposition                        RETIRE
old delivery-targeted constraints
  RETIRE old form unless successor authored
D-3.3 placement proposition
  REPLACE when successor meaning is materialized
D-3.4/D-3.5 related elements
  STALE FOR REVIEW where relevance was transfer-dependent
\end{verbatim}

\textbf{Finding C:} no successor is required merely because a
neighboring function remains. \texttt{RETIRE} and \texttt{REPLACE} are
materially different.

\subsection{6. Test D - order-fulfillment internal versus external
WMS}\label{test-d---order-fulfillment-internal-versus-external-wms}

The order corpus states that \texttt{ADR-2.1} internal inventory
authority is a genuine responsibility-boundary choice: an external WMS
would replace many project-owned inventory FRs with
integration/API/contract obligations.

This provides an independent domain control.

Expected BA effect:

\begin{itemize}
\tightlist
\item
  stable order/evaluation and governed result meanings may retain
  identity;
\item
  internal inventory authority/capability/store meanings may retire or
  be replaced;
\item
  internal mutation propositions retire/supersede;
\item
  external service/contract referents and consumption/normalization
  propositions are introduced;
\item
  payment and fulfillment meanings outside the affected boundary do not
  become stale merely because MR-2 changed.
\end{itemize}

\textbf{Finding D:} mutation-local lifecycle handling generalizes beyond
the camera corpus without requiring a domain taxonomy.

\subsection{7. Family-specific identity
tests}\label{family-specific-identity-tests}

\subsubsection{7.1 BAReferent test}\label{bareferent-test}

Retain identity when the independently reusable project-semantic thing
remains the same despite changed propositions around it.

Mutation examples forcing retention:

\begin{itemize}
\tightlist
\item
  \texttt{LocalConnectivity} across Ethernet/Wi-Fi;
\item
  \texttt{RecognitionCapture} across remote/local recognition when
  capture itself remains required;
\item
  order-evaluation/result semantics across provider-boundary changes
  where governed meaning is preserved.
\end{itemize}

\subsubsection{7.2 BAProposition test}\label{baproposition-test}

Retain identity only under normalized assertion equivalence.

A proposition identity therefore cannot survive a material change in:

\begin{itemize}
\tightlist
\item
  operator;
\item
  polarity;
\item
  role-bound participant identity/value;
\item
  governed constraint value;
\item
  local condition/temporal scope.
\end{itemize}

Source wording and provenance may change without forcing replacement.

\textbf{Result:} one shared retention rule for both BA1 families is
rejected.

\subsection{8. Review-state pressure}\label{review-state-pressure}

Origin state alone cannot encode review state.

Examples:

\begin{verbatim}
DERIVED + ACCEPTED
DERIVED + REJECTED
DIAGNOSTIC_UNRESOLVED + ACCEPTED
GROUNDED + STALE relative to new baseline
\end{verbatim}

Therefore
\texttt{PENDING\_REVIEW\ \textbar{}\ ACCEPTED\ \textbar{}\ REJECTED} is
required as a separate analytical review dimension.

This does not give BA project authority. It records whether analytical
materialization is admitted to the accepted Base Analysis.

\subsection{9. Staleness pressure}\label{staleness-pressure}

A new governed baseline does not invalidate all old BA meaning.

The minimum staleness rule is dependency-local:

\begin{verbatim}
changed source locator
  -> directly grounded dependent BA elements need review

stale/replaced/retired derivation basis
  -> derived dependents need review

replaced/retired proposition participant
  -> proposition needs review
\end{verbatim}

M3 also shows that effective context can matter even where the direct
source file of an element did not change.

Therefore exact change-impact closure remains open, but \texttt{STALE}
is required as a target-baseline revalidation state.

\subsection{10. Supersession versus
retirement}\label{supersession-versus-retirement}

The mutation corpus forces both:

\subsubsection{SUPERSEDED}\label{superseded}

A successor exists and takes over the earlier semantic responsibility.

Examples:

\begin{itemize}
\tightlist
\item
  Ethernet realization proposition -\textgreater{} Wi-Fi realization
  proposition;
\item
  negative ownership proposition -\textgreater{} affirmative ownership
  proposition;
\item
  remote placement proposition -\textgreater{} local placement
  proposition.
\end{itemize}

\subsubsection{RETIRED}\label{retired}

The old meaning ceases and no successor identity is required.

Examples:

\begin{itemize}
\tightlist
\item
  remote capture-transfer proposition after local recognition;
\item
  delivery-behavior referent if no longer used;
\item
  external service-consumption proposition after project ownership, if
  no external consumption remains.
\end{itemize}

Collapsing both into \texttt{inactive} would lose change explanation and
targeted downstream invalidation.

\subsection{11. Diagnostic resolution
test}\label{diagnostic-resolution-test}

Suppose baseline B0 contains an accepted \texttt{DIAGNOSTIC\_UNRESOLVED}
element because sources conflict or are insufficient. B1 corrects the
governed documentation.

Rejected behavior:

\begin{verbatim}
same diagnostic identity
originState: DIAGNOSTIC_UNRESOLVED -> GROUNDED
\end{verbatim}

That rewrites history.

Accepted provisional behavior:

\begin{verbatim}
issue disappears       -> diagnostic RETIRE
new unresolved issue   -> diagnostic REPLACE / SUPERSEDE
new factual meaning    -> separate GROUNDED/DERIVED BA element
                         with own provenance
\end{verbatim}

Thus origin remains historical while lifecycle expresses resolution.

\subsection{\texorpdfstring{12. Why BA-level \texttt{REVISE} is not
required}{12. Why BA-level REVISE is not required}}\label{why-ba-level-revise-is-not-required}

The source corpus itself uses review labels such as \texttt{revise},
\texttt{revise/replace} and \texttt{retire/supersede} for governed
documents.

Those labels are not a BA identity ontology.

At BA semantic identity level:

\begin{itemize}
\tightlist
\item
  non-identity metadata/provenance refresh with same meaning
  -\textgreater{} \texttt{RETAIN};
\item
  changed meaning with successor -\textgreater{} \texttt{REPLACE};
\item
  disappeared meaning -\textgreater{} \texttt{RETIRE}.
\end{itemize}

Adding \texttt{REVISE} would blur the exact boundary T2 is trying to
protect: whether semantic identity survived.

\subsection{13. Minimum candidate after
pressure}\label{minimum-candidate-after-pressure}

\begin{verbatim}
Origin/provenance dimension:
  GROUNDED | DERIVED | DIAGNOSTIC_UNRESOLVED

Review dimension:
  PENDING_REVIEW | ACCEPTED | REJECTED

Freshness dimension:
  CURRENT | STALE

Cross-baseline continuity disposition:
  RETAIN | REPLACE | RETIRE

Derived historical interpretation:
  accepted REPLACE -> SUPERSEDED
  accepted RETIRE  -> RETIRED
\end{verbatim}

No new BAE family is forced.

\subsection{14. Negative controls}\label{negative-controls}

\begin{itemize}
\tightlist
\item
  Do not make Git commit ancestry the semantic identity rule.
\item
  Do not use document IDs as BA proposition identity.
\item
  Do not preserve proposition identity solely because source wording is
  similar.
\item
  Do not replace a referent merely because
  ownership/provider/realization changed.
\item
  Do not auto-transfer retired security constraints to a new behavior.
\item
  Do not mark the whole BA stale because one source file changed.
\item
  Do not treat \texttt{REJECTED} analytical material as governed project
  truth.
\item
  Do not mutate a resolved diagnostic into a grounded fact.
\end{itemize}

\subsection{15. BA1 / BA2 reopen check}\label{ba1-ba2-reopen-check}

No mutation requires a third identity family or a new
operator/role/modifier semantic.

The pressure is absorbed by:

\begin{itemize}
\tightlist
\item
  BA1 identities;
\item
  BA2 proposition structure;
\item
  BA3 provenance + continuity/lifecycle metadata.
\end{itemize}

\begin{verbatim}
BA1 reopen trigger: false
BA2 reopen trigger: false
\end{verbatim}

\subsection{16. Exit disposition}\label{exit-disposition}

\begin{verbatim}
BA3-T2              COMPLETED / PROVISIONAL PASS
                    WITH IDENTITY-LIFECYCLE REFINEMENT
BA3                 STARTED / NOT CLOSED
BA1                 REMAINS CLOSED
BA2                 REMAINS CLOSED
\end{verbatim}

The remaining BA3 work is narrower: derivation-rule reproducibility,
effective dependency/change-impact lineage and feedback provenance,
followed by closure review.

\subsection{17. Next falsification
target}\label{next-falsification-target}

Only after this package is committed, pushed and remotely verified:

\begin{quote}
\textbf{BA3-T3 - derivation-rule reproducibility and change-impact
lineage pressure test.}
\end{quote}

T3 should use the facial M4 effective-governed-context pressure plus
provider-state normalization/order boundary cases. It must not start BA4
or method-specific analysis schemas.

\end{document}
